\section{Introduction}
\label{sec:introduction}

The explosive advancement and widespread deployment of Large Language Models (LLMs) have fundamentally reshaped digital communication, automated content synthesis, and information dissemination. State-of-the-art autoregressive foundation models generate synthetic text of unprecedented fluency, often indistinguishable from human writing across diverse registers~\citep{touvron2023llama,solaiman2019release}. While these generative capabilities unlock profound societal benefits in education, automated translation, and productivity, they simultaneously present acute epistemic hazards. The proliferation of automated synthetic text enables the scalable propagation of coordinated disinformation campaigns, algorithmic astroturfing in consumer reviews, automated academic dishonesty, and the degradation of open-web knowledge repositories with plausible yet entirely fabricated claims~\citep{chakraborty2023possibilities,sadasivan2023can}.

In response, the Natural Language Processing (NLP) community has developed diverse methodologies for AI-generated text detection, ranging from zero-shot statistical curvature estimators~\citep{mitchell2023detectgpt,bao2024fastdetectgpt,hans2024binoculars} and token rank rankers~\citep{gehrmann2019gltr} to supervised deep transformer discriminators~\citep{solaiman2019release,wang2024raid}. However, the overwhelming majority of existing detection frameworks and empirical benchmarks have been developed and evaluated almost exclusively on high-resource, morphologically isolating or fusional languages, predominantly English. Low-resource and typologically diverse languages---particularly agglutinative languages of the Turkic family spoken by over 180 million people worldwide---remain largely unaddressed in the generative AI auditing literature.

Among these languages, \textbf{Kazakh}---a Kipchak Turkic tongue spoken by over twenty million individuals across Central Asia---presents an extraordinary combination of structural and sociolinguistic properties that expose severe systemic failure modes in standard neural text processing paradigms:
\begin{enumerate}
    \item \textbf{Agglutinative Morphotactic Complexity}: Kazakh grammar operates through extensive linear affixation chains, wherein multiple inflectional and derivational suffixes sequentially attach to nominal and verbal roots to express plurality, possession, case relations, tense, aspect, mood, voice, negation, and evidentiality~\citep{makhambetov2013kazakh,bekmanova2021kazakh}. A single verbal stem such as \textit{жаз} (to write) gives rise to hundreds of distinct surface realizations (e.g., \textit{жазбағандықтан} ``because [someone] did not write'', \textit{жаздырылмағандықтан} ``because [it] was not caused to be written''). Standard subword tokenization algorithms, such as Byte-Pair Encoding (BPE)~\citep{sennrich2016subword} and SentencePiece~\citep{kudo2018sentencepiece}, are predominantly parameterized on massive Indo-European corpora. When applied to Kazakh, these tokenizers partition inflected word forms into arbitrary, semantically fragmented byte fragments. This phenomenon, which we term \textit{subword token inflation}, destroys morpheme boundaries, exponentially inflates sequence lengths, and induces severe out-of-vocabulary sparsity in pretrained language models.
    \item \textbf{The False Positive Crisis on Short and Vernacular Texts}: In high-stakes real-world auditing environments (such as educational grading, consumer protection, and content moderation), the primary operational hazard is the rate of \textit{false positive accusations}---falsely identifying authentic human writing as machine-generated. False positives inflict severe academic, legal, and reputational harm on human authors~\citep{liang2023gptdetectors}. We observe that this failure mode is dramatically exacerbated on short texts ($\le 60$ characters), such as consumer product reviews, customer support messages, and social media comments, where subword token fragmentation deprives classifiers of contextual clues, leading to catastrophic false-alarm rates.
    \item \textbf{Multi-Domain Shift and Bilingual Code-Switching}: Real-world Kazakh digital ecosystems, particularly on mobile platforms (Telegram, Instagram, VKontakte), are characterized by pervasive bilingual code-switching between Kazakh and Russian, featuring hybrid morphosyntactic loanwords (e.g., Russian commercial stems bearing Kazakh case inflections, such as \textit{доставкасы} ``its delivery'', \textit{каспиден} ``from Kaspi''). Standard AI detectors trained on clean formal prose experience severe degradation when deployed across such colloquial out-of-distribution domains.
    \item \textbf{The Orthogonal Imperative of Factual Veracity}: A critical blind spot of contemporary AI text detection is that detecting machine generation does not answer whether a statement is \textit{factually true}. Conversely, verifying factual claims without examining generative provenance leaves systems defenseless against synthetic hallucinations. An authoritative encyclopedia article written by a human may contain factual errors, while a generative model may produce rigorously grounded, factually accurate text. Isolated single-axis classifiers cannot untangle these dimensions, exposing automated verification pipelines to severe manipulation.
\end{enumerate}

\subsection{Research Questions}
To establish a rigorous, principled foundation for generative AI detection and factual verification in agglutinative low-resource languages, this investigation addresses three central research questions:
\begin{itemize}
    \item \textbf{RQ1 (Morphological Grounding)}: \textit{To what extent does explicit finite-state morphological segmentation and morpho-syntactic diagnostic feature fusion mitigate subword token inflation and eliminate false positive accusations on short texts compared to standard subword representations?}
    \item \textbf{RQ2 (Cross-Domain Generalization and Code-Switching Robustness)}: \textit{How robust are monolingual versus multilingual neural architectures when exposed to severe domain shifts, colloquial social media registers, and bilingual code-switched loanwords, and can multi-generator contrastive representation learning prevent generator overfitting?}
    \item \textbf{RQ3 (Joint Credibility Calibration)}: \textit{How can factual veracity and machine-generation likelihood be integrated into a unified continuous multidimensional calibration space to reliably separate authentic knowledge, human misconceptions, AI hallucinations, and verified AI-assisted writing?}
\end{itemize}

\subsection{Principal Contributions}
In addressing these questions, this paper establishes a comprehensive, end-to-end framework uniting morphological linguistic theory, multi-generator contrastive learning, hybrid retrieval, and multidimensional trust modeling. Our principal contributions are fivefold:
\begin{enumerate}
    \item \textbf{Linguistic Formalization and Finite-State Transducer (FST)}: We formally articulate the morphotactic rules governing Kazakh nominal cases, verbal tense/aspect chains, negation allomorphs, and evidentials. We develop an 83-rule Finite-State Transducer (FST) that explicitly decouples native inflectional affixes and bilingual code-switched loanword stems, proving mathematically that FST segmentation curbs subword token inflation by up to 36.65\% and forces subword tokenizers to converge to the empirical Kazakh morphemic limit of 2.202 tokens per word.
    \item \textbf{Comprehensive Multi-Benchmark Corpus Suite}: We construct and release three extensive, human-validated evaluation benchmarks for Kazakh NLP:
    \begin{itemize}
        \item \textbf{KazAI-Detect}: A balanced multi-domain benchmark comprising 12,848 in-domain consumer review pairs across Appstore, Bookstore, Mapping, and Market categories, complemented by a 1,317-sample Out-of-Distribution (OOD) test suite spanning news and Wikipedia.
        \item \textbf{Kaz-MAGE / Kaz-RAID}: A multi-generator adversarial benchmark capturing social media comments across Telegram, VKontakte, and TikTok with natural Russian-Kazakh code-switching.
        \item \textbf{Kazakh-FEVER 3K}: The first human-validated fact verification benchmark for Kazakh, featuring 3,000 claim-evidence pairs with an explicit \textsc{Hard NEI} protocol targeting morphological negation, evidential, and temporal traps.
    \end{itemize}
    \item \textbf{Unified Morphologically-Grounded Neural Framework}: We design a dual-stream architecture that dynamically fuses dense pretrained transformer representations with explicit morphological boundary signals via gated cross-attention, regularized by a supervised contrastive multi-generator objective ($\mathcal{L}_{\text{total}} = \mathcal{L}_{\text{CE}} + \lambda \mathcal{L}_{\text{SupCon}}$) that prevents overfitting to generator-specific decoding artifacts.
    \item \textbf{16-Dimensional Morphological Diagnostic Cross-Encoder and 2D Trust Matrix}: We introduce a 16-dimensional morphological diagnostic vector capturing verbal polarity XOR ($\Delta_{\text{neg}}$), calendar year divergence ($\Delta_{\text{year}}$), evidential past markers (\textit{-ыпты/-іпті}), and FST root Jaccard overlap, effectively eliminating lexical shortcut traps. We unify this verifier with generative detection into a continuous Cartesian 2D Trust Matrix ($\mathcal{T}(x) = \sqrt{\frac{1}{2}(\max(0, T_{\text{fact}})^2 + T_{\text{gen}}^2)}$), enabling calibrated risk assessment.
    \item \textbf{Empirical Validation and Open-Source Deployment}: Extensive GPU experiments across four model families (mBERT, XLM-RoBERTa, KazBERT, KazRoBERTa) demonstrate that our morphological framework slashes false positive accusations on short texts by 43.21\% ($p < 0.00001$), improves factual verification Macro-F1 by 7.8\%, and maintains over 98.8\% accuracy on out-of-distribution news and Wikipedia. We deliver an end-to-end interactive demonstration system featuring real-time SVG Cartesian plane rendering and expandable morphological diagnostic drawers deployed on Hugging Face Spaces.
\end{enumerate}

The remainder of this article is structured as follows. Section~\ref{sec:related_work} synthesizes related work across statistical AI detection, low-resource Turkic NLP, and automated fact verification. Section~\ref{sec:linguistic_foundation} details the linguistic foundations of Kazakh morphotactics, the 83-rule FST analyzer, and subword segmentation pathologies. Section~\ref{sec:framework_architecture} presents the unified architecture encompassing gated fusion, contrastive learning, hybrid retrieval, and the 2D Trust Matrix. Section~\ref{sec:experimental_benchmarks} reports extensive empirical benchmarks and ablation findings. Section~\ref{sec:system_and_explainability} describes the production system and interactive explainability interface. Finally, Section~\ref{sec:conclusion_and_future_work} concludes with ethical considerations, limitations, and future research trajectories.
